\subsection{Post-saturation novelty refresh after claim-expansion search}

A later claim-expansion search on 2026-08-19 surfaced additional direct pressure on the interpretation of P1.D2. Preregistered Belief Revision Contracts (PBRC) already supplies an operator-agnostic protocol layer in which admissible belief change is tied to preregistered evidence triggers, allowed revision operators, priorities, fallback, and externally validated evidence witnesses \cite{pbrc2026}. Paper~I therefore does not claim generic evidence-gated belief revision or a generic control layer above revision operators. Self-Revising Discovery Systems for Science already treats scientific discovery as a verified transition between representational regimes with typed artifacts, provenance, preservation and transport obligations \cite{selfrevisingdiscovery2026}; and the Model Discovery Agent performs M-open hypothesis-class expansion after predictive inadequacy and chooses informative experiments to identify the proposed mechanism \cite{modeldiscoveryagent2026}. Generic scientific regime revision, model-class expansion, and active discrimination are therefore donor territory rather than Paper~I novelty.

The same refresh also tightens the failure-to-repair boundary. Model or Harness? formulates an interaction-centric repair-assignment problem in which the same visible failure can belong to the model, harness, environment, or grader and should be repaired at the responsible component \cite{modelorharness2026}. Diagnosis Is Not Prescription further shows that the component most causally responsible for an agent failure need not be the best intervention target \cite{diagnosisnotprescription2026}. Paper~I therefore does not claim generic component-responsibility localization, repair assignment, or the diagnosis-versus-intervention distinction. Finally, Earned Authority under a Fixed Ceiling for Evolving Agents formalizes authorization continuity across agent mutation and allows evidence-gated authority changes below an immutable effect ceiling \cite{earnedauthority2026}; generic mutation authorization and non-amplification are not Paper~I novelty either.

These sources narrow the novelty boundary, while the later faithful-comparator challenge narrows the empirical boundary. The phrase ``authority gate'' is not a standalone novelty claim. The v2.2.4 archive shows how the governed policy combines lower-level exclusion, same-task repair, protected preservation, and dependency-impact binding. It does not show that this composition is comparatively necessary. An active value-of-information parent with ordered candidate search matches the governed policy on the primary joint criterion. The earlier margin over non-iterating parents is therefore retained as diagnostic history rather than comparative mechanism evidence.

The subsequent P1-X study tests a wider but still bounded systems residual: whether the distinction among causal responsibility, best intervention, and scientifically justified revision matters after the donor mechanisms above are already granted. Its comparator-fair disjoint V2 supports that joint escalation contract on five exact heterogeneous revision-responsibility families, while the ideal information-equivalent product ties exactly (Section~\ref{sec:p1x}). Accordingly, the novelty claim is the explicit \emph{composition and interface discipline} of minimal scientific escalation under the tested contracts, not ownership of belief revision, diagnosis, repair assignment, model expansion, authorization, or any inherent centralization/expressivity advantage. Open-ended or deployed scientific-agent generality remains undetermined.
